\documentclass[10pt,a4paper]{amsart}
\usepackage[margin=19mm]{geometry}
\usepackage{amsmath,amssymb,mathtools}
\usepackage[scr=rsfs,cal=euler]{mathalfa}
\usepackage{xfrac}
\usepackage[unicode,hidelinks,bookmarks=false]{hyperref}
\newcommand{\ie}{i.e.\ }
\newcommand{\cf}{cf.\ }
\newcommand{\resp}{respectively\ }
\newcommand{\Subsection}{\subsection}
\providecommand{\nobreakdash}{\nobreak\hbox{-}}
\setlength{\emergencystretch}{2em}
\multlinegap0pt
\usepackage{amssymb}
\usepackage{dsfont}
\usepackage{xcolor}
\usepackage[normalem]{ulem}
\usepackage{enumerate}
\usepackage{tikz}
\usepackage{mathtools}
\usepackage[shortlabels]{enumitem}
\newtheorem{lemma}{Lemma}
\newcommand{\N}{\mathbb{N}}
\title{A Banach space with an unconditional basis which is not slicely countably determined}
\author{Marcus L\~oo, Yo\"el Perreau}
\date{Source: arXiv:2603.12947v1 (CC BY 4.0)}
\begin{document}
\maketitle

\noindent\emph{Source and licence.} A Banach space with an unconditional basis which is not slicely countably determined. Version arXiv:2603.12947v1 (CC BY 4.0), distributed by arXiv under the Creative Commons Attribution 4.0 International licence, http://creativecommons.org/licenses/by/4.0/, as recorded in the arXiv licence metadata for this exact version and shown on the abstract page https://arxiv.org/abs/2603.12947v1. Full licence notice: \texttt{licenses/arxiv-2603.12947-CC-BY-4.0.txt}.\par\smallskip

\noindent\textit{Editorial scope.} Counted unit: the article's lemma lemma:functinals\_on\_X\_infty\_are\_essentially\_supported\_on\_a\_finitely\_branching\_tree (source0.tex lines 560--562). The article's complete original proof of it is delivered with it (source0.tex lines 564--585, one \texttt{proof} environment of 22 lines). No mathematics is written, repaired or re-derived: the statement and the proof are the article's own lines, in the article's own notation, and no retained line is substituted or re-flowed.  No further source-local context is needed: every cross-reference of every retained line resolves inside the two delivered spans.  Nothing was omitted from any retained span, so the package carries no omission record.  No retained line contains a citation, so the wrapper carries no bibliography and no bibliography entry.  No retained line contains a figure environment, an image-inclusion command or drawing code, so no image is delivered and the package declares no asset dependency.  Editorial formatting only, in addition: an AMS \texttt{amsart} preamble that restates the article's own declarations and macros used by the retained lines, namely the environments \texttt{lemma} on separate counters, together with the standard packages those lines need.  The retained environments print in this document as Lemma 1 in order of appearance, whereas the article prints the same items with the numbering of its own full document.  The complete native source of the article as distributed in the arXiv e-print for this exact version is bundled unchanged as \texttt{sources/196304/source0.tex}.
\begin{lemma}\label{lemma:functinals_on_X_infty_are_essentially_supported_on_a_finitely_branching_tree}
    Let $f\in X_{T_\infty}^*$ and $\varepsilon>0$. There exists a finitely branching tree $T_0\subseteq T_{\infty}$ such that $||P^{*}_{T_{\infty}\setminus T_0}f||<\varepsilon$.
\end{lemma}

\begin{proof}
 Fix a non-zero $f\in X_{T_\infty}^*$ and $\varepsilon>0$. For convenience, given two subsets $A,B$ of $T_\infty$, we will write $A\preceq B$ if every element of $B$ is a successor of some element of $A$.  First, we construct inductively a sequence $(R_k)_{k\in \N\cup\{0\}}$ of finite subsets of $T_\infty\cup\{\emptyset\}$ with $R_0=\{\emptyset\}$ and such that for every $k\in\mathbb{N}$, the following conditions are met.
    \begin{enumerate}
        \item $R_k\subseteq L_k$;
        \item $R_k\succeq R_{k-1}$;
        \item $||P^*_{S_k}f||<\varepsilon/2^k$, where $S_k = \bigcup_{s\in L_k\setminus R_k}T(s)$.
    \end{enumerate} Indeed, let $k=0$ or assume that for some $k\in\N$ we have found $R_k$ satisfying the conditions (1), (2) and (3). Observe that \begin{equation*}
        ||P^*_{\bigcup_{\substack{s\in L_{k+1}\\
    s\succeq R_k}}T(s)}f|| = \sum_{\substack{s\in L_{k+1}\\ s\succeq R_k}} ||P^*_{T(s)}f||,
    \end{equation*}
    so we can find a finite subset $R_{k+1}\succeq R_k$ of $L_{k+1}$ such that\begin{equation*}
        \sum_{s\in R_{k+1}}||P^*_{T(s)}f|| >  ||P^*_{\bigcup_{\substack{s\in L_{k+1}\\ s \succeq R_{k}}}T(s)}f||- \frac{\varepsilon}{2^{k+1}}.
    \end{equation*}
    In particular, \begin{equation*}
        ||P^*_{S_{k+1}}f|| = ||P^*_{\bigcup_{s\in L_{k+1}\setminus R_{k+1}}T(s)}f|| = \sum_{s\in L_{k+1}\setminus R_{k+1}}||P^*_{T(s)}f||<\frac{\varepsilon}{2^{k+1}},
    \end{equation*} so $R_{k+1}$ satisfies all the required conditions. Finally, let $T_0:= \bigcup_{k\in\mathbb{N}} R_k$. Then $T_0$ is a finitely branching subtree of $T_\infty$, and 
\begin{equation*}
    T_{\infty}\setminus T_0 = \bigcup_{k\in\mathbb{N}}\bigcup_{s\in L_k\setminus R_k} T(s) = \bigcup_{k\in\mathbb{N}}S_k.
\end{equation*} Since the $S_k$ are totally incomparable, it follows that \begin{equation*}
     || P^*_{T_{\infty}\setminus T_0}f|| = ||P^*_{\bigcup_{k\in\mathbb{N}}S_k}f|| = \sum_{k\in\mathbb{N}}||P^*_{S_k}f|| < \sum_{k\in\mathbb{N}}\frac{\varepsilon}{2^k}=\varepsilon,
\end{equation*} as we wanted.
\end{proof}

\end{document}
